%% How EXHALE solves a new composition from a converged wind (restart path),
%% and what the lower-atmosphere handoff carries: the scalar base.inp against
%% the profile file. Both parts record implemented behavior read off the
%% source; Part 2 also keeps the interface decision that led to the profile
%% route and checks the plan's requirement list against what was built.
%% Updated 2026-09-10: the restart contract (Restart intent, the eight-line
%% metadata block and its refusals, Restart option change, the grid guard,
%% the provenance-unknown mark), the element counts with the oxygen-chemistry
%% species, the He/H comparison over the physical cells and the message that
%% names the deviating cell, and the mass-conserving element projection.
%% Author: Kwang-Il Seon
\documentclass[english,a4paper]{my_memo}
\synctex=1
\usepackage{booktabs}
\usepackage{amsmath}
\usepackage{graphicx}
\usepackage{babel}
\emergencystretch=3em

\title{Solving a new composition from a converged wind, and\\
       what the base handoff carries}
\author{Kwang-Il Seon}
\date{Last updated: \docmoddate}

\begin{document}
\maketitle

\section{What this memo records}

Two things that are easy to confuse with each other.

\medskip
\noindent\textbf{Part 1: the composition restart.} How a converged EXHALE
solution is used as the starting guess for a \emph{different} composition:
what the run script does, what \verb|load_IC| does to the loaded densities,
where the seed's information is kept and where it is overwritten, and
(the point that matters for how results are quoted) why the seed sets only
the path and never the answer. Everything in this part is implemented,
verified against the source, and in production use in the LHS\,1140\,b
composition and $K_{zz}$ scans.

\medskip
\noindent\textbf{Part 2: the base handoff.} EXHALE now has two handoff
routes: the scalar \path{base.inp} and the profile file named by
\verb|Lower atmosphere profile:|. This part records what each one is in the
code today, the comparison on which the choice between them was made, and
(item by item against the plan's requirement list) what the profile
route actually carries, what reads each field, and what is still not
transmitted.

\medskip
\noindent\textbf{Decision (2026-08-26), implemented 2026-08-27.} Phase\,E
goes to a profile file. The scalar \path{base.inp} is not retired: it stays
as the single-layer route, and its behavior with no profile beside it is
unchanged. What the decision left to the implementation (the schema, and
what happens when both files are present) is now written, and a scalar
physics key beside a profile is refused rather than merged
(\path{src/modules/files_IO/input_read.f90:1371-1386}).

\bigskip
\hrule

\part*{Part 1: solving a new composition from a converged wind}
\addcontentsline{toc}{part}{Part 1}

\section{The recipe, as the run script executes it}
\label{sec:recipe}

The production example is \path{LHS1140b/exhale/run_heh_diff_case.sh}, which
solves one point of the helium-abundance scan with the binary H/He diffusion
operator active. In order (line numbers are that file):

\begin{enumerate}
\item \textbf{Clone the input, change one line.} The control case's
  \verb|input.inp| is copied with \verb|sed| and exactly the composition
  line is replaced (lines 27--31):
  \verb|He/H number ratio:| takes the new value, while the planet tag, the
  direct-steady switch and the residual target are set for this run and
  \verb|He_Kzz:| is appended (line 32). Nothing else about the
  configuration moves (same spectrum, same domain, same flags), so the
  scan is a scan in composition and not in anything else.
\item \textbf{Hand the converged seed over as the initial condition}
  (lines 33--34): the seed run's \verb|output/Hydro_ioniz.txt| and
  \verb|output/Ion_species.txt| are copied into the new run directory as
  \verb|output/Hydro_ioniz_IC.txt| and \verb|output/Ion_species_IC.txt|,
  which are the two file names \verb|load_IC| reads
  (\path{src/modules/files_IO/load_IC.f90:2-4,89,106,126}).
\item \textbf{Solve the steady state directly.} The binary is run with
  \verb|EXHALE_PTC=1|, \verb|EXHALE_PTC_JFNK=1| and
  \verb|EXHALE_PTC_DTAU0=1.0| (line 43), i.e.\ the JFNK direct-steady route
  rather than a marching relaxation.
\item \textbf{Feed the output back as the next initial condition}, when the
  diffusion operator is on (lines 40--41): each further invocation copies
  the run's own last output over its \verb|*_IC.txt| files. This is the
  outer Picard loop of composition and wind driven from the outside, one
  pass-set at a time, and the loop is stopped when the reported composition
  drift stops falling.
\item \textbf{Post-processing and transit} (lines 51--57): one
  \verb|Do only PP: True| pass to write the advection-corrected profiles,
  then \verb|EXHALE_transit.py| for the He\,\textsc{i} 10830 line.
\end{enumerate}

Inside a single invocation the same Picard loop runs up to twenty times when
\verb|He_diffusion| is on and exactly once when it is off
(\path{src/EXHALE_main.f90:1238}); it exits on a composition drift below
$10^{-3}$ (line 1268) or on a failed steady solve (line 1254), and the
composition update is under-relaxed with $\omega$ starting at 0.5 and halved
down to a floor of 0.125 on any pass whose drift failed to fall (lines
1229, 1270--1274). The drift tested is the undamped distance to the fixed
point, so a small $\omega$ cannot buy a false convergence (lines
1213--1214).

\section{What the restart does to the composition}
\label{sec:redistribute}

The seed file holds \emph{absolute} species densities. Read as they stand,
they would silently run the donor's composition while the setup report
echoed the input one. The input file is the authority, so \verb|load_IC|
carries the loaded hydrogen and helium onto the requested ratio
(\path{src/modules/files_IO/load_IC.f90:428-443}).

\subsection{Counting the elements}

In each cell the two element densities are counted over every species that
carries a nucleus, the molecular carriers and the oxygen-chemistry species
included (\path{load_IC.f90:444-452}):
%
\begin{align}
  n_{\rm H} &= n(\mathrm{H\,\textsc{i}}) + n(\mathrm{H\,\textsc{ii}})
     + 2\,[\,n(\mathrm{H_2}) + n(\mathrm{H_2^+})\,]
     + 3\,n(\mathrm{H_3^+}) + n(\mathrm{HeH^+})
     + n(\mathrm{OH}) + 2\,n(\mathrm{H_2O}),
  \label{eq:nH}\\
  n_{\rm He} &= n(\mathrm{He\,\textsc{i}}) + n(\mathrm{He\,\textsc{ii}})
     + n(\mathrm{He\,\textsc{iii}}) + n(\mathrm{HeH^+}).
  \label{eq:nHe}
\end{align}
%
$\mathrm{HeH^+}$ is inside both counts, because it carries a nucleus of
each. The metastable $2^3S$ state is \emph{not} a term of its own: the
\verb|HeI| column the code transports and the file carries is the total
He\,\textsc{i} population, the triplet included, so its nucleus is already
counted there and adding it again would count it twice
(\path{load_IC.f90:449-450}).

These are the integer weights \verb|bsp_nH| and \verb|bsp_nHe| of the
species table (\path{src/modules/init/species_table.f90:155-160}), over the
thirteen base species
\verb|HI HII HeI HeII HeIII HeTR H2 H2+ H3+ HeH+ OH H2O CO|:
%
\begin{center}\small
\begin{tabular}{@{}lccccccccccccc@{}}
\toprule
 & HI & HII & HeI & HeII & HeIII & HeTR & H$_2$ & H$_2^+$ & H$_3^+$
 & HeH$^+$ & OH & H$_2$O & CO \\
\midrule
\verb|bsp_nH|  & 1 & 1 & 0 & 0 & 0 & 0 & 2 & 2 & 3 & 1 & 1 & 2 & 0 \\
\verb|bsp_nHe| & 0 & 0 & 1 & 1 & 1 & 1 & 0 & 0 & 0 & 1 & 0 & 0 & 0 \\
\bottomrule
\end{tabular}
\end{center}
%
taken together with \verb|bsp_is_excited_level|, \verb|.true.| for
\verb|HeTR| alone, which is what every budget sum skips
(\path{species_table.f90:185-189}). The same count is what
\verb|element_ratio_HeH| applies in the code proper
(\path{src/modules/functions/composition.f90:374-402}) and what the
diffusion operator writes back against
(\verb|element_nucleus_counts|, called from \verb|mixture_mass_split|,
\path{src/modules/functions/binary_element_diffusion.f90:3041-3106}). One
element count, one definition, three call sites: there is no second
formula that can drift away from this one. That the triplet had to be
excluded was measured, not assumed: it was being added a second time in all
three places and in four more, and the mass, particle, nucleus and
collision-partner sums were each wrong by that amount until 2026-08-27
(\path{docs/Update_EXHALE_stage1.pdf} \S75).

\subsection{The redistribution}

Let $\mathrm{HeH}$ be the requested ratio. The deviation of the loaded state
is measured cell by cell as $|n_{\rm He}/n_{\rm H} - \mathrm{HeH}|/
\mathrm{HeH}$ and the redistribution is triggered when its maximum exceeds
\verb|heh_dev_tol| $=10^{-6}$ (\path{load_IC.f90:479-493}). A restart at the
composition the file was written with therefore leaves every density
untouched (bit for bit), which is what keeps the same-composition
restart path exact.

\paragraph{The comparison runs over the physical cells alone, and it names
the cell it attained.} Two properties of that loop were measured into it and
both matter to anyone reading a refusal message. First, the range is
$j = 1 \ldots N$: the ghosts are boundary data, the lower pair being the
inflow reservoir whose molecular partition the handoff imposes and the sweep
re-pins every step, and the upper pair mirroring the top cell, so a check on
the input's He/H has nothing to say about them. The reason it is not the
whole array is a measured incident (2026-09-05): a carrier-transport
snapshot carried its lower ghost with He/H $2.1\times10^{-5}$ off the input
while every physical cell agreed to $10^{-6}$, and that ghost alone sent the
restart into the rescale branch, which the $\mathrm{HeH^+}$ clause below
then refused. Second, the quantity that decides is the \emph{largest}
departure over those cells, and \verb|report_deviating_heh_cell| prints the
cell that attained it, with its radius, the He/H the file carries there, the
He/H the input asks for, the relative departure and the threshold, all at
\verb|ES16.9|. The message used to print the top physical cell instead,
which on a molecular state is not the deciding cell at all: MEASURED on a
coupled-carrier state, the top cell agreed with the input to seventeen
digits while cell 199 at $r = 1.0906$ was off by $3.96\times10^{-2}$
relative, so the refusal read ``7.9307E-02 against 7.9307E-02 \ldots{}
cannot be rescaled''. Two numbers that print the same cannot explain a
refusal.

When it does trigger, the two scale factors are
(\path{load_IC.f90:514-515})
%
\begin{equation}
  s_{\rm H} = \frac{n_{\rm H} + n_{\rm He}}
                   {(1 + \mathrm{HeH})\, n_{\rm H}},
  \qquad
  s_{\rm He} = \mathrm{HeH}\,
               \frac{n_{\rm H} + n_{\rm He}}
                    {(1 + \mathrm{HeH})\, n_{\rm He}},
  \label{eq:scales}
\end{equation}
%
so that the new element densities in that cell are
%
\begin{equation}
  n'_{\rm H} = \frac{n_{\rm H} + n_{\rm He}}{1 + \mathrm{HeH}},
  \qquad
  n'_{\rm He} = \mathrm{HeH}\,\frac{n_{\rm H} + n_{\rm He}}
                                   {1 + \mathrm{HeH}},
  \qquad
  n'_{\rm H} + n'_{\rm He} = n_{\rm H} + n_{\rm He}.
  \label{eq:conserve}
\end{equation}
%
That is the defining property of the operation: \textbf{the total nuclei
count of each cell is conserved} and only the split between the two elements
moves. Every hydrogen-bearing column is multiplied by $s_{\rm H}$ and every
helium-bearing one by $s_{\rm He}$
(\path{load_IC.f90:530-547}), the trace metals following hydrogen because
they are defined relative to a hydrogen nucleus, and OH, H$_2$O and CO
following it too because among the two elements being rescaled they carry
only H (their oxygen and carbon follow hydrogen through \verb|melem_ab|,
like every trace metal), so
%
\begin{itemize}
\item the \textbf{ionization-stage split inside each element} is preserved
  exactly: $n(\mathrm{H\,\textsc{ii}})/n_{\rm H}$,
  $n(\mathrm{He\,\textsc{ii}})/n_{\rm He}$ and the $2^3S$ fraction all come
  through unchanged, since a single factor multiplies every stage of the
  element;
\item the \textbf{velocity, pressure and temperature} profiles are the
  seed's, read straight off \verb|Hydro_ioniz_IC.txt|;
\item the \textbf{mass density is not} read from the file. It is rebuilt by
  \verb|calc_rho| from the redistributed densities under the run's own mass
  policy, He\,$2^3S$, molecular and (under \verb|eos_metals|) metal mass
  included (\path{load_IC.f90:778-800}). It therefore does move with the
  composition, as it must: swapping hydrogen for helium at fixed nuclei
  count changes the mass. The seed supplies the kinematics and the thermal
  state; the mass follows the new composition.
\end{itemize}

\subsection{Two refusals, and the diffusion-on exception}

$\mathrm{HeH^+}$ carries one nucleus of each element, so no single factor
can set both counts. Two cases are therefore refused outright rather than
approximated (\path{load_IC.f90:494-513}): a file carrying no helium at all
(nothing to rescale), and a molecular file carrying $\mathrm{HeH^+}$ whose
composition disagrees with the input while \verb|He_diffusion| is off:
such a state must be restarted at its own composition.

With \verb|He_diffusion| on the rule changes, because the cell-by-cell
element split \emph{is} the state being restarted: the diffusion operator
produced it, and a separated profile is the physics, not a defect of the
file. The rescale is then confined to the base cells $1-N_g \ldots 1$,
which is exactly where the operator itself holds a Dirichlet reservoir
composition; every cell above keeps the helium fraction it was written with,
$s_{\rm H} = s_{\rm He} = 1$ (\path{load_IC.f90:516-529}). In those base
cells $\mathrm{HeH^+}$ is projected onto the two element totals the same way
\verb|binary_element_diffusion| writes back: it takes the smaller of the two
factors, $\min(s_{\rm H}, s_{\rm He})$, and whatever nuclei that leaves
short of $n'_{\rm H}$ and $n'_{\rm He}$ are deposited into the neutral
ground species $\mathrm{H\,\textsc{i}}$ and $\mathrm{He\,\textsc{i}}$
(\path{load_IC.f90:548-576}; the same projection at
\verb|project_elements|, \path{binary_element_diffusion.f90:2107-2266}).
The projection is
non-negative by construction (the smaller factor can only reduce
$\mathrm{HeH^+}$), and it conserves both element totals, at the cost of
moving a little of the molecular ion into the neutral species.

Either way the run says on stdout which of the two it did: with
\verb|He_diffusion| on it prints the kept top-cell ratio and the input
ratio, and otherwise it prints the input ratio followed by the deviating
cell's line (\path{load_IC.f90:574-585}), so the restart is never silent.

\subsection{Elements the file does not carry}

Separately from the composition, an element whose stages the file does not
carry (no column at all, or columns identically zero because the run that
wrote them had metals off) is built from its abundance exactly as a cold
start does, $n_X = \mathrm{ab}_X\, n_{\rm H}$ with everything in the neutral
stage, and the substitution is reported
(\path{load_IC.f90:589-648}). The column test alone would accept a
metals-off file and then run with zero metal density while the base boundary
condition still counted metals in the mass and particle budget, leaving an
$O(1)$ residual in the first cells; the zero-sum test is what closes that.

\subsection{Elements whose reservoir the handoff has moved}
\label{sec:handoff-rescale}

A third case sits between the two above and was added on 2026-08-27: the
file carries the element, but the \emph{reservoir} has moved since the file
was written. That is exactly what a flux-closure iteration does (it
changes the elemental ratios of the lower atmosphere from one iteration to
the next), so the base boundary condition would use the new $\mathrm{El/H}$
while every cell above the base still held the old one, and the elemental
budget $n_{\rm El}/n_{\rm H} = (\mathrm{El/H})$ could not close above the
first cell (\path{docs/Update_EXHALE_stage1.pdf} \S79--80).

The rule is a single factor on the whole column
(\path{load_IC.f90:650-705}):
%
\begin{equation}
  r_{\rm El} = \frac{(\mathrm{El/H})_{\rm handoff}}
                    {(\mathrm{El/H})_{\rm loaded,\ base\ cell}},
  \label{eq:rel}
\end{equation}
%
applied to every ionization stage of that element alike, so the loaded
ionization split and the radial shape of the profile cancel out of it and
only the normalization moves. $\mathrm{El/H}$ at the base cell is counted in
nuclei (the element summed over its stages against the hydrogen nuclei of
equation~\eqref{eq:nH}), which is the convention of
\verb|element_ratio_HeH| and of \path{src/utils/element_budget.py}, so the
factor is measured in the units the budget check reads it back in.

Which elements are touched is not a test on the numbers. An element is
rescaled if and only if the lower-atmosphere handoff \emph{states} its
reservoir, and what records that is the door the value came in through:
\verb|set_element_abundance| is the only routine either handoff reader uses
to set an abundance, and passing through it sets
\verb|melem_from_handoff(e)| (\verb|set_element_abundance| in
\path{src/modules/files_IO/input_read.f90}). An abundance that
reached the run from \path{metals.inp} alone never passes through that door,
so it is not marked and its column is loaded untouched; a restart with no
handoff at all leaves the loop doing nothing at all, which is what keeps
every pre-existing restart bit for bit what it was. Three further guards
skip the element: one already rebuilt from the abundance by the block above
(it \emph{is} the reservoir already), a zero abundance, and a loaded column
that sums to zero at the base cell.

Helium is deliberately outside this rule and keeps the older convention of
sections~\ref{sec:redistribute}: the base cells are set to the input
$\mathrm{He/H}$ and, with \verb|He_diffusion| on, the column above keeps its
diffused split.

Each rescaled element prints its own line: loaded ratio, stated ratio,
factor, at enough digits that a factor of $1.0001$ is legible, because that
is the size of the step a converged closure iteration takes
(\path{load_IC.f90:701-704}).

\section{The restart contract: what a state file states, and what is refused}
\label{sec:contract}

Everything above is what the loader does to the numbers. This section is
what the loader checks before it does any of it. A state is a solution of an
equation set on a grid at a composition, so a file and a run that do not
agree on those three things are not a restart pair at all, and since
2026-09-09 the disagreement is named rather than absorbed.

\subsection{What a restart is for: \texttt{Restart intent}}

A loaded state can be continued in three senses, and they are not variants
of one path: a physical trajectory carries the clock the file records and
every step obeys the physical-mode contract; a relaxation carries no clock
and marches toward stationarity; a stationary evaluation takes no step at
all and measures the state as it stands. One selector says which,
%
\begin{center}
\verb+Restart intent: trajectory | relaxation | stationary [evaluate|equilibrate]+
\end{center}
%
because a run that does not state its purpose cannot be checked against it
(\path{src/modules/files_IO/input_read.f90}). The second word belongs to
\verb|stationary| alone: \verb|evaluate| returns the state unchanged with
its certification and takes no solver step, \verb|equilibrate| puts the
loaded composition on its own fixed point before the evaluation, which is a
change of the state and is reported when it is taken, and the bare
\verb|stationary| measures the state as loaded and then enters the
stationary solve.

Four consistency checks refuse rather than reinterpret, because a key that
is silently ignored is a run doing something other than what its input file
says: the key with \verb|Load IC? False| (there is no loaded state for the
intent to be about), \verb|Restart option change| with the same
(no state is being restarted), \verb|trajectory| with \verb|Run mode: init|
(a trajectory is a physical integration and an initialization run claims no
elapsed time), and \verb|stationary| without \verb|Solver: Newton| (there is
no stationary solve to enter). The default is the run mode's own meaning:
\verb|trajectory| in \verb|phys|, \verb|relaxation| otherwise, which is what
every restart did before the key existed, so no existing input changes
behavior.

\verb|stationary_state_of_the_loaded_restart| (\path{src/EXHALE_main.f90})
is the route that measures the state as loaded. Every earlier path into the
solver went through the marching loop, which takes at least two CFL steps
before it hands off, and a step is a different state, so what used to be
certified was never the state the file carried. This route takes no step: it
fixes WENO3, rebuilds the derived quantities in the checkpoint order
(composition, then boundary, then primitive state), runs one equilibrium
sweep because that sweep is part of the residual's definition, prints the
loaded composition's departure from that sweep's root, and holds the
residual and the certification to the file's own \verb|certified=| field.
MEASURED: the certified \verb|wasp_full_newton| state is not the fixed point
of its own sweep (the particle count moves $2.1\times10^{-11}$, the species
$5.0\times10^{-10}$, the cooling $3.6\times10^{-10}$, which is the cell
solve's $x_{\rm tol} = \sqrt{\epsilon}$ and the file's digits), and that is
reported at every evaluation rather than equilibrated away.

\subsection{The metadata block, and the six fields it compares}

Both state files carry an eight-line block of \verb|#| comments, built and
parsed by one routine each (\verb|build_restart_metadata| and
\verb|verify_restart_metadata|, \path{src/modules/files_IO/load_IC.f90}),
with every number written \verb|ES23.16| so that a field compared as text is
compared as the number it stands for. Every reader of these files skips
\verb|#| lines, so the block changes no numeric parse and no regression
verdict.

\begin{center}\small
\begin{tabular}{@{}lll@{}}
\toprule
Field & What it states & On disagreement \\
\midrule
\verb|restart_schema| & the block's own version & refuse \\
\verb|reservoir| & He/H and every trace El/H the run carries & refuse, naming the element \\
\verb|species_columns| & the columns the state carries & report \\
\verb|grid| & $N$, $R_0$, $r(1)$, $r(N)$, the grid mode & refuse \\
\verb|constants| & the constant set, $R_J$, $k_B$ & refuse \\
\verb|options| & 20 physics switches, token by token & refuse, naming the token \\
\verb|t_phys[s]| & the clock of the state & adopted if the coupling line lacks it \\
\verb|source| & the build that produced it & report \\
\bottomrule
\end{tabular}
\end{center}

The division is the point. The \verb|options| field carries every switch
that adds or removes an equation, an unknown or a source term, so a
difference means the file's state solves another system; the numeric inputs
(fluxes, tolerances, rates) are deliberately \emph{not} there, because they
change the coefficients of the same system and a state carried across such a
change is a legitimate starting point whose residual simply is not zero. The
\verb|species_columns| line is reported and never refused, because species
are restored by the label on the \verb|# columns| line and an element the
file does not carry is built from its abundance (\S\ref{sec:redistribute}).
A pair whose two halves carry different blocks is refused outright: they
were not written by one run, so they are not one state.

\paragraph{A file with no block is loaded, and marked.} A state written
before the block existed is loaded exactly as it was, and the run says so:
the configuration it was produced under cannot be compared, so every file
that run writes carries \verb|restart_input=provenance_unknown| in its
\verb|source| field. The mark is a property of the state and is inherited,
because a state descended from an unknown configuration is no better known
at the second reload than at the first, and dropping the mark there would
lose it for good.

\paragraph{The grid belongs to the run.} \verb|verify_restart_radii|
compares the cell centers of the physical cells $1 \ldots N$ to $10^{-10}$
relative and refuses any other grid, naming the worst cell and both radii.
The reason is that \verb|define_grid| builds the faces, the widths and the
window indices $j_{\min}$, $j_{\rm flux}$ together with the centers, and a
state file carries only the centers, so a file written under a different
\verb|Grid type|, \verb|Base grid| or \verb|Outer radius| passes a row-count
guard whenever it has the same \verb|Grid cells| and would then be marched
on one grid's centers with another grid's faces. The tolerance is chosen
against two measurements: the smallest real grid difference in the tree is
$3.3\times10^{-4}$, and the coarsest writer of a state file, the Wind-AE
initial condition at \verb|ES17.10|, has a worst relative round-off of
$5.0\times10^{-11}$, half of the tolerance. This reader does not
interpolate.

\subsection{Naming what a restart may change: \texttt{Restart option change}}

The exact \verb|options| comparison forbids the workflow this project
actually uses to reach its solutions, which is to converge without an option
and restart with it on. The key
\verb|Restart option change: <token>[, <token> ...]| names the tokens that
are allowed to differ; everything else still refuses by name, with the
token and its \verb|from -> to|. Six layout tokens (\verb|metals|,
\verb|mol|, \verb|oxychem|, \verb|carrier|, \verb|carrier_newton|,
\verb|iontrans|) may never be named, because they decide which rows the
state carries and naming one cannot make the file's rows this run's rows:
that is a cold start, not a restart. An unknown token, a grid, reservoir or
constant word, the key with no token, and the key without
\verb|Load IC? True| are all input errors.

A permitted change is written into both halves of the new state as
\verb|# option_change from -> to at restart of <source>| and inherited, up
to a history of 32 lines, and \verb|EXHALE_setup.out| gains a
\verb|Restart provenance:| line reading \verb|UNKNOWN|,
\verb|UNKNOWN by descent| or \verb|states its configuration| together with
the changes this run permits.

\paragraph{What the state files no longer report.} Two fields left the
coupling header when the base boundary became the characteristic face
condition: \verb|valve| and \verb|fluxconst| were the two switches the old
ghost-cell boundary chose for itself, and the present boundary has neither,
so a header reporting them would be reporting settings the run cannot have.
The loader's parser reads the fields it recognizes one at a time and ignores
the rest, so a file written before that change still restores its
\verb|sec_ion| state and an older reader still finds \verb|sec_ion| in a
file written after it.

\section{The seed is a guess; the answer is the new steady state}
\label{sec:seed}

This is the distinction the whole recipe rests on, and it is worth stating
plainly because the two are easy to conflate when a scan is quoted as a
table of restarts.

\medskip
\noindent The redistribution of section~\ref{sec:redistribute} is
\emph{arithmetic on a starting guess}. It is not a solution at the new
composition, and it is not claimed to be one: it holds $(v, p, T)$ and every
ionization-stage ratio at the donor's values, which are the values that were
in equilibrium with the \emph{donor's} composition. What makes the result a
result is what happens afterwards: the JFNK steady solve, the
ionization-equilibrium solve, and (with diffusion on) the composition outer
loop, iterated until the wind residual meets its target and the composition
drift stops moving. The convergence \emph{path} is what the seed choice
buys; the converged state is a property of the new configuration alone.

\medskip
\noindent The operational consequence: a restart chain is legitimate exactly
as long as each rung is converged on its own terms, and each row of a scan
must be quoted with its own achieved residual and drift, not with the
seed's. That is why the LHS\,1140\,b tables carry \verb|info|, $\|R\|$ and
the composition drift in every row.

\medskip
\noindent Two measured cases show that the seed choice is a real question of
path and nothing more.

\paragraph{The helium scan needed no ladder.} In the $K_{zz} = 10^{9}$
composition scan, every case restarted from the converged
\emph{diffusion-off} solution of the same composition, or from the
neighboring composition already solved with diffusion on, and every case
reached \verb|info = 0| \emph{in its first pass-set} and then stopped moving
(\path{LHS1140b/kzz_decision.md}, section~6). The recorded reason is
physical: a diffusion-off solution is well mixed near the base, which is
where the adopted $K_{zz}$ also mixes, so the step the line search is asked
to take is small. The same held at $10^{8}$ and $10^{10}$
(\emph{ibid.}, section~6.1).

\paragraph{The eddy coefficient did need one.} Restarting a case at
$K_{zz} = 10^{8}$ or above directly from the $K_{zz} = 0$ control asks the
solver to absorb the whole helium column in one step: the line search stalls
and the composition outer loop stops on its solver-failure guard after a
single pass (\emph{ibid.}, section~3). Rebuilding one decade at a time
(\path{ladder_kzz.sh} restarts each case from the case one decade below it)
fixes the seeding. It does not fix everything: on this wind JFNK settles
at $\|R\|$ between $2.7\times10^{-3}$ and $4.7\times10^{-3}$ rather than the
$1.0\times10^{-3}$ the control reaches, so the upper rungs state a looser
\verb|Resid tol| and quote the achieved $\|R\|$ in the table. A target of
$2\times10^{-2}$ was tried and rejected, because there the solver stops
responding to the eddy term altogether.

\medskip
\noindent Same code, same operation, opposite verdicts on whether a ladder
is required, which is the concrete demonstration that seeding is about
the step the nonlinear solver is asked to take, and that no property of the
converged answer follows from it.

\section{Evidence that the restart is faithful}
\label{sec:evidence}

\subsection{The element projection conserves mass}
\label{sec:projection-mass}

Every write of an element total back into a species vector goes through
\verb|project_elements| (\path{binary_element_diffusion.f90}), and what it
has to preserve is the mass the species handed in carried. Until 2026-09-10
it took the mixture mass as
$m_{\rm sum} = m_1 n_{\rm H} + m_{\rm He} n_{\rm He}$ with $m_1$ the
\emph{reservoir} mass per hydrogen nucleus, which carries the trace metals
at \verb|melem_ab| while the cell carries them at whatever the trace
transport left, and the difference was written into hydrogen at every call.
MEASURED at one cell of the certification column, with metal/H $0.024401$
against the reservoir's $0.028227$: $\sum_i m_i f_i$ went
$1.0000000 \to 1.0000636$ after one step, helium's gain exceeding the
hydrogen and metal loss by $6.4\times10^{-5}$, which is mass created. A
second leak of the same kind had the trace-metal transport writing metal densities
with no counter-flux in the hydrogen they diffuse through
($3.0\times10^{-3}$ on the suite column with the metals frozen).

\verb|mixture_mass_split| now reads the mass the cell's species actually
carry, from \verb|bsp_mass| and \verb|melem_A| with no species weight
restated, and \verb|project_elements| scales component~1 as a whole, every
helium-free species included, with the metals written first and the hydrogen
group closing the mixture; one projection after the metal loop puts the
metals' mass change back on hydrogen. Metals off is bitwise by construction
($m_1 = m_{\rm H} = 1$ exactly). MEASURED: the certification column's
closure after the relaxation went $3.37\times10^{-3} \to 2.3\times10^{-15}$,
and the handed-back atomic candidate $6.44\times10^{-3} \to
1.94\times10^{-9}$, which is the file's own precision. Reported and not
closed: the metal transport's counter-flux is closed on hydrogen but not on
momentum or internal energy, which is a trace approximation.

\subsection{The state survives the file}

\paragraph{The writer-to-loader round trip returns the state.} MEASURED on
the 12000-step \verb|mol_base_handoff| state with \verb|EXHALE_DUMP_IC=1|,
which writes the state as loaded: $r$, $v$ and $p$ come back bitwise, every
species column and every inventory (H, He, O and C nuclei, charge, He/H, the
H$_2$ ratio) to $4.2\times10^{-16}$, $T$ to $5.6\times10^{-16}$, and a second
trip to $5.0\times10^{-16}$ with no drift. The \verb|rho| column comes back
at $4.9\times10^{-12}$, and that is the writer's own mass closure and not the
loader's: the loader never reads that column, it rebuilds the mass from the
species, and perturbing the column by $10^{-15}$ changes nothing bitwise
(the same state at 40 steps closes to $4.1\times10^{-14}$).

The stationary residual of a reloaded state was a separate matter and was
closed later. Before the element projection was fixed to conserve mass
(\S\ref{sec:projection-mass}), the writer and the reload of one atomic
candidate stood at $\|R\| = 3.2\times10^{-2}$ and $0.4385$ \emph{on
different rows}; after it, at $1.112$ and $1.1123$ on the \emph{same} row.
The state is the same object on both sides of the file.

\paragraph{What the restart does \emph{not} fix, and where the movement
comes from.} MEASURED, and worth separating from the loader: He/H departs
from the input by $1.0\times10^{-11}$ and $8.9\times10^{-12}$ after 12000
marching steps with carrier transport off and on, but by 3.9 to 7.9 percent
on a coupled-carrier density-unknown JFNK state and by 41 to 43 percent on
an $\ln n$ state. The carrier column is written after the element projection
precisely so that the Newton's carrier is not rescaled, the sweep reads the
nuclei totals from the state, and with helium diffusion off no row re-imposes
He/H. That is the element constraint of the stationary system seen from the
restart side; it is the producing path that moved the ratio, not the loader.

\paragraph{The diffused profile survives a round trip.} Measured against the
stored \path{examples/14_diffusion} output under \verb|EXHALE_DUMP_IC=1|,
which writes the state as loaded, before the first ionization-equilibrium
solve: the loaded He/H profile reproduces the file's to
$3.3\times10^{-16}$ relative, worst cell, over all 504 cells
(\path{docs/Update_EXHALE_stage1.pdf} \S67). That is round-off, i.e.\ the
\verb|He_diffusion| branch of section~\ref{sec:redistribute} leaves the
column it is meant to leave alone.

\paragraph{The defect it replaced.} The same file through the pre-change
binary comes back \emph{flattened} to the input ratio: at $r = 1.05\,R_p$
the file holds $1.616\times10^{-2}$ and the load produced
$8.333\times10^{-2}$, a factor 5.2 (\emph{ibid.}). The old rule rescaled
every cell to the global \verb|HeH| whenever the file's composition differed
from it, on the ground that the input file is the authority: correct with
diffusion off, and destructive of the answer with diffusion on. Both
numbers are quoted from that record and were not re-measured for this memo.

\paragraph{And the composition authority itself was a measured need.} A
He/H\,$=10$ restart seeded from a He/H\,$=1$ solution was measured to
converge back to He/H\,$=1$ everywhere before the redistribution existed
(the measurement is recorded at the head of that block in
\path{load_IC.f90}); the donor's composition ran while the setup report
echoed the input one.

\paragraph{The diffusion-off path is unmoved.} \verb|make check| reported
byte-identical output on every case of the matrix as it stood on 2026-08-25,
which is the evidence that confining the rescale to the base cells changed
nothing outside the \verb|He_diffusion| branch
(\path{docs/Update_EXHALE_stage1.pdf} \S67, quoted; the matrix has since grown
from five cases to sixteen, and was not re-run for this memo).

\bigskip
\hrule

\part*{Part 2: the base handoff (scalars and profiles)}
\addcontentsline{toc}{part}{Part 2}

\section{The scalar handoff}
\label{sec:scalars}

\verb|read_base_inp| reads an optional \path{base.inp} from the run
directory (\path{src/modules/files_IO/input_read.f90:1203-1388}). It
recognizes six named keys and one family of element keys, and every key
belongs to one of the categories the routine's own header lists, because the
category is what says what the value is allowed to do to the wind:

\begin{center}
\begin{tabular}{@{}lllp{5.6cm}@{}}
\toprule
key & unit & category & what it sets \\
\midrule
\verb|T_base|    & K       & EOS boundary &
   \verb|T0|, the base temperature \\
\verb|r_base|    & $R_J$   & EOS boundary &
   \verb|R0|, the base radius \\
\verb|p_base|    & bar     & EOS boundary &
   \verb|p_base_bar|, the level all of the above refer to
   (default $10^{-6}$) \\
\verb|q_H2_base| & ratio   & EOS boundary &
   \verb|q_h2_base|, the H$_2$ volume mixing ratio that sets the base
   particle count through \verb|comp_ntot_bc|, in place of the
   chemical-equilibrium fit. An EOS anchor, not a composition pin: nothing
   holds H$_2$ there and no H$_2$ profile is seeded from it \\
\verb|HeH_base|  & ratio   & elemental reservoir &
   \verb|HeH|; a positive value also sets \verb|thereis_He| \\
\verb|<El>_H_base| & ratio & elemental reservoir &
   $X_{\rm El}$, the El/H nuclei ratio, through
   \verb|set_element_abundance|, for any of the ten elements of the species
   table (e.g.\ \verb|O_H_base 4.90e-4|). Overrides \path{metals.inp} for
   that element, and switches the metal system on if it is the only nonzero
   abundance \\
\verb|Kzz_base|  & cm$^2$\,s$^{-1}$ & boundary constraint &
   \verb|he_kzz|, the eddy coefficient the element-diffusion operator
   imposes at the base; inert with \verb|He_diffusion| off \\
\bottomrule
\end{tabular}
\end{center}

Comment lines and unknown keys are ignored, with one comment parsed:
\verb|# solution_id <hash>|, which is what the pair check of
section~\ref{sec:profile} reads. The call sits in \verb|input_read| after
\verb|run_lower_atm_prestep| and \emph{before} the derived constants, so the
overrides take effect on everything computed from them, and an absent file
is a no-op (\path{input_read.f90:858-882}). Two consistency messages are
printed and not enforced: that \verb|q_H2_base| and \verb|HeH_base| must
come from the same lower-atmosphere solution
(\path{input_read.f90:1325-1336}), and that a \verb|p_base| away from
1\,$\mu$bar means the base composition and the equilibrium H$_2$ fit both
refer to that other level (\path{input_read.f90:1337-1342}). Those two stay
messages deliberately: an internally inconsistent scalar handoff is the
user's to judge. What is \emph{not} left to a message is the pair of files,
see section~\ref{sec:profile}.

\verb|Kzz_base| deserves a note, because the scalar it writes has changed
underneath it. \verb|he_kzz| now defaults to $0$, not $10^{9}$: an eddy term
is a property of the atmosphere being modeled, so it is stated rather than
inherited (\path{src/modules/init/parameters.f90:146-156}). And
\verb|he_kzz| is no longer what the diffusion operator reads. Every site in
\verb|binary_element_diffusion| reads the array \verb|kzz_cell| instead
(\path{parameters.f90:157-164}; the three sites are
\path{binary_element_diffusion.f90:679,1380,1597}), filled once the radial
grid exists by \verb|eddy_diffusion_on_grid|: from the interpolated
\verb|Kzz| column of a profile when there is one, and from \verb|he_kzz| in
every cell otherwise
(\path{src/modules/files_IO/lower_atmosphere_profile.f90:465-529}). The
uniform case is exactly the old scalar (for a uniform array the face
average $\tfrac12(a+a)$ is $a$ in IEEE double), so a run with no profile
sees the number it always saw.

A real file, the one the \verb|mol_base_handoff| regression case pins
(\path{backup/regression/mol_base_handoff/base.inp}, comments elided):

\begin{quote}\ttfamily
T\_base\ \ \ \ 1140.0\\
r\_base\ \ \ \ 0.49\\
HeH\_base\ \ 0.0793\\
Kzz\_base\ \ 1.000e+09\\
q\_H2\_base 0.75\\
p\_base\ \ \ \ 1.000e-05
\end{quote}

Four of the seven default regression cases carry such a file:
\verb|mol_base_handoff|, \verb|mol_metals|, \verb|mol_lyman_werner| and
\verb|mol_diffusion| (\path{backup/regression/run_check.sh:57}), so the
scalar contract is pinned by the golden comparison and cannot be changed
silently. Those four goldens were re-snapshotted on 2026-08-27, but for a
reason that has nothing to do with the handoff: a molecular-network audit
found the three-body term using a mass density where a number density
belongs (\path{docs/Update_EXHALE_stage1.pdf} \S82). The scalar contract itself is unchanged.

\section{The profile handoff}
\label{sec:profile}

The second route is the file named by \verb|Lower atmosphere profile:| in
\path{input.inp} (\path{input_read.f90:447-455}; user manual \S3.8). It is a
table over an interval of pressure rather than a single level, and the
reader is \path{src/modules/files_IO/lower_atmosphere_profile.f90}.

\paragraph{What the file is.} A header of \verb|# key value| lines and then
one row per level, deep to shallow. The header states the fingerprint
(\verb|solution_id|, \verb|source_code|, \verb|source_version|,
\verb|mechanism|, \verb|stellar_flux|, \verb|notes|), the three pressures
(\verb|p_match_bar|, \verb|p_top_bar|, \verb|p_deep_bar|), the closure
iteration (\verb|iteration|, \verb|trial_flux_H|, \verb|trial_flux_He|,
\verb|reached_steady_state|) and the column names
(\path{lower_atmosphere_profile.f90:185-198}). Nine columns are required
(\verb|p|, \verb|r|, \verb|T|, \verb|n_tot|, \verb|rho|, \verb|Kzz|,
\verb|q_H2|, \verb|q_H|, \verb|X_He|), and the element columns
\verb|X_<El>| are found by name, so the producers may order them as they
like (\path{lower_atmosphere_profile.f90:328-340}).

\paragraph{What is refused.} The reader refuses rather than warns on a
non-positive or non-monotonic pressure, a table whose top does not reach
above the matching level (\verb|p_top_bar| $\ge$ \verb|p_match_bar|, i.e.\
no overlap at all), a matching level outside the pressures the table covers,
and a missing \verb|solution_id|
(\path{lower_atmosphere_profile.f90:344-382}). Interpolation is linear in
$\log p$ and is never extrapolated; a matching level that falls on a node
returns that node bit for bit, which is what lets a single-level profile
reproduce the scalar handoff exactly
(\path{lower_atmosphere_profile.f90:396-434}).

\paragraph{What EXHALE takes from it.} Exactly five kinds of quantity, all
through the same doors the scalar keys use
(\verb|apply_lower_atmosphere_profile|, \path{input_read.f90:1392-1460}):
\verb|T0| and \verb|R0| from the \verb|T| and \verb|r| columns at the match,
\verb|p_base_bar| from \verb|p_match_bar|, \verb|q_h2_base| from the
\verb|q_H2| column, \verb|HeH| from \verb|X_He|, and each metal abundance
from its \verb|X_<El>| column through \verb|set_element_abundance|, the
same routine \verb|<El>_H_base| goes through, which is why an element stated
by either route is marked \verb|melem_from_handoff| and picked up by the
restart rescale of section~\ref{sec:handoff-rescale}. \verb|K_zz| is taken
as a profile, interpolated onto the radial grid into \verb|kzz_cell|. The
\verb|n_tot| and \verb|rho| columns are required of the file but are
\emph{not} imposed: the base density is \verb|n0| of \path{input.inp}, and
the base particle count and mass follow from $T_0$, He/H, $q_{\rm H_2}$ and
the elemental ratios exactly as on the scalar path.

\paragraph{The two files together are enforced, not printed.} With a
profile in use, every \path{base.inp} key in the three physics categories of
section~\ref{sec:scalars} (EOS boundary, elemental reservoir, boundary
constraint) is refused with \verb|error stop|, because the profile
already states that quantity at the matching level and a scalar accepted
beside it would be a second source of the same number
(\verb|refuse_scalar_key|, \path{input_read.f90:1371-1386}). What a
\path{base.inp} beside a profile may still carry is provenance, and there
the pair is checked: a missing \verb|# solution_id| line, or one that
differs from the profile's, stops the run
(\path{input_read.f90:1344-1367}). \verb|He_Kzz:| in \path{input.inp} is a
different file with a different contract and is not refused, but it is inert
beside a profile and the run says so
(\path{input_read.f90:1449-1458}).

\paragraph{What pins it.} The \verb|lower_profile| regression case, one of
the seven defaults (\path{backup/regression/lower_profile/}): HD 209458 b
with \verb|He_diffusion: True|, the base state and $q_{\rm H_2}$ taken from
the table at $p_{\rm match} = 10^{-6}$ bar, the elemental reservoirs
carried by the profile with \emph{no} \path{metals.inp} in the directory
(so the profile is the only source of the metal abundances), a
non-constant \verb|Kzz| column interpolated onto the grid in place of the
inert scalar, and a provenance-only \path{base.inp} whose \verb|solution_id|
matches, which pins the accepting branch of the pair check. The worked
example is \path{examples/17_lower_profile/}.

\section{Scalar file against profile file}
\label{sec:compare}

This is the comparison the choice was made on, brought up to what each route
is in the code today.

\begin{center}
\small
\begin{tabular}{@{}p{2.4cm}p{5.3cm}p{5.7cm}@{}}
\toprule
 & scalar \path{base.inp} & profile file \\
\midrule
what is transmitted
 & six named numbers plus one El/H per element, all describing one layer
 & columns on the lower model's own pressure grid over an interval that
   reaches above the match \\[2pt]
relation between the models
 & a handoff: the lower model's state at one level becomes the upper
   model's boundary
 & a matching over an interval both models cover, refused outright when the
   file does not reach above the match \\[2pt]
flux continuity
 & not checkable. One layer's values carry no gradient, so nothing in the
   file constrains an elemental flux, and \verb|HeH_base| stays a
   prescribed input while ``the wind returns H:He'' is an assumption
   dressed as a result
 & checkable, and measured. The elemental face flux
   $4\pi r^2 (\rho X v + J)$ is reduced to a median and a relative spread
   over two radial windows and written to \path{EXHALE_resolved.out}
   whenever a profile is in use
   (\path{binary_element_diffusion.f90:1784-1872}); the criterion is the
   same radial invariance as T8 (\path{docs/Update_EXHALE_stage1.pdf} \S69) \\[2pt]
eddy coefficient
 & one constant, \verb|Kzz_base| $\to$ \verb|he_kzz|
 & the \verb|Kzz| column, interpolated onto every cell as \verb|kzz_cell|
   \\[2pt]
molecular composition
 & one scalar, \verb|q_H2_base|
 & \verb|q_H2| and \verb|q_H| at their own levels, plus the diagnostic
   molecules; the ions of the EXHALE network are still not transmitted
   (section~\ref{sec:carries}) \\[2pt]
what it costs to produce
 & a keyword scan, unknown keys ignored
 & a schema with units and a grid convention, interpolation onto the EXHALE
   grid, a fingerprint and a stated failure policy, all written
   (\path{src/utils/lower_profile_schema.py} and the two adapters beside
   it) \\[2pt]
failure mode
 & a silent mismatch: the two consistency checks are messages, not refusals
   (\path{input_read.f90:1325-1342}), so a handoff assembled from two
   different lower-atmosphere solutions runs without complaint, unless a
   profile is present, in which case the pair is checked by fingerprint
 & a stated one: a missing or unmatched \verb|solution_id|, a scalar
   physics key beside the profile, a table that does not overlap the match,
   each an \verb|error stop| with the reason printed \\
\bottomrule
\end{tabular}
\end{center}

Three rows carried the decision.

\paragraph{Flux continuity is the whole point of Phase\,E.} A steady
lower-atmosphere chemistry solution cannot determine the hydrogen supply
independently of its upper-boundary flux. Without a flux-continuity
condition at a matching surface, \verb|HeH_base| remains prescribed and the
composition is an input wearing the clothes of an output. A scalar file
cannot express the condition, because a single level has no gradient and
carries no flux. This is not a convenience argument; it is the difference
between a coupling and an assumption.

\paragraph{A constant $K_{zz}$ is the wrong shape.} The judgment is already
on record: ``a \emph{constant} $K_{zz}$ is the wrong shape of assumption for
this problem regardless of its magnitude: the molecular coefficient it
competes against rises by nearly five decades between the base and
$1.2\,R_p$, so no single constant can both mix the base and follow the
molecular coefficient outward, while Lindzen's own result is that the eddy
coefficient should \emph{turn over} rather than keep growing''
(\path{docs/eddy_diffusion_kzz.tex}, summary). \verb|Kzz_base| can transmit
only a constant. The homopause (the one consequential thing $K_{zz}$
does) is set where $K_{zz} = D$, so it is a property of the two profiles
and not of either number alone.

\paragraph{The molecular scalar has the same defect one level down.}
\verb|q_H2_base| is a single H$_2$ mixing ratio standing in for a
composition that the molecular network resolves into H$_2$, H$_2^+$,
H$_3^+$ and HeH$^+$. It is enough to select a branch of the base particle
count; it is not enough to hand over a photochemical solution. This is the
one row of the three the profile route did not close: it transmits
$q_{\rm H_2}$ and $q_{\rm H}$ at their own levels, but none of the three
molecular ions (section~\ref{sec:carries}).

\section{The decision, and what was built on it}
\label{sec:decision}

\subsection{Decided 2026-08-26 (user decision)}

\begin{enumerate}
\item \textbf{Phase\,E goes to a profile file.} The interface decision the
  plan left open (``whether this stays a scalar file or becomes a profile
  file is an interface decision to confirm with the user before
  implementation'') is settled in favor of profiles, on the flux-continuity
  row of section~\ref{sec:compare} above all. The plan itself now records
  the phase as done (milestones E1--E5, 2026-08-27) with its three
  acceptance criteria and the numbers behind them.
\item \textbf{The scalar \path{base.inp} is kept}, as the single-layer
  route rather than a deprecated one. Three reasons, in order: four
  regression cases pin it and their goldens are the evidence that the
  molecular base is unmoved; a single-level handoff remains the right
  description when the lower model genuinely produces one (the analytic
  column of Tier\,1); and the code path is small and already validated.
\item \textbf{The profile route is opt-in through a new key}
  (\verb|Lower atmosphere profile:|), and its
  absence leaves the previous behavior exactly as it was. This is the
  standing house rule: a new physics option defaults to off so old
  goldens still reproduce.
\end{enumerate}

\subsection{What was implemented, 2026-08-27}

All three hold, with one qualification on the third. The route is opt-in
through \verb|Lower atmosphere profile:| and an absent key is a no-op, so
the diffusion-off and scalar-handoff paths are untouched by any of it. The
qualification is that the four molecular goldens are \emph{not} the ones
they were at the time of the decision: they were re-snapshotted at the end
of the same day for the three-body density fix of
\path{docs/Update_EXHALE_stage1.pdf} \S82, which is a different change with a
different reason. The Phase\,E work did not move them.

What the decision left open, and where it now lives:

\begin{center}
\small
\begin{tabular}{@{}p{4.0cm}p{9.4cm}@{}}
\toprule
piece & where it is \\
\midrule
the schema
 & \path{docs/input_schema.md} \S2d; user manual \S3.8 \\
the reader
 & \path{src/modules/files_IO/lower_atmosphere_profile.f90} \\
the run key
 & \verb|Lower atmosphere profile:| (\path{input_read.f90:447-455}) \\
the producers
 & \path{src/utils/lower_profile_schema.py} (format, elemental accounting,
   fingerprint) with \path{photochem_to_lower_profile.py} and
   \path{vulcan_to_lower_profile.py} behind one command line \\
the closure loop
 & \path{src/utils/element_flux_closure.py}, driven from the fluxes EXHALE
   measures and writes to \path{EXHALE_resolved.out} \\
the regression
 & \path{backup/regression/lower_profile/}, one of the seven defaults \\
the example
 & \path{examples/17_lower_profile/} \\
\bottomrule
\end{tabular}
\end{center}

\subsection{What the profile file carries, against what the plan asked for}
\label{sec:carries}

The plan's list (\emph{ibid.}, Phase\,E part 1) against the schema, field by
field. ``read by'' names the code that consumes it; a field the file carries
and nothing consumes is stated as such rather than counted as met.

\begin{center}
\small
\begin{tabular}{@{}p{3.5cm}p{3.4cm}p{6.5cm}@{}}
\toprule
asked for & field & read by \\
\midrule
matching pressure and radius
 & header \verb|p_match_bar|; columns \verb|p|, \verb|r|
 & \verb|p_base_bar| directly; \verb|R0| from \verb|r| at the match
   (\path{input_read.f90:1418-1424}) \\[2pt]
temperature and total density
 & columns \verb|T|, \verb|n_tot|, \verb|rho|
 & \verb|T0| from \verb|T|. \verb|n_tot| and \verb|rho| are required of the
   file but not imposed: the base density stays \verb|n0| of
   \path{input.inp} and the base particle count and mass are built from
   $T_0$, He/H, $q_{\rm H_2}$ and the elemental ratios, as on the scalar
   path \\[2pt]
elemental H and He abundances
 & columns \verb|X_He|, \verb|X_<El>|
 & \verb|HeH| and, through \verb|set_element_abundance|, every
   \verb|melem_ab| (\path{input_read.f90:1430-1443}) \\[2pt]
\textbf{the upward elemental fluxes}
 & header \verb|trial_flux_H|, \verb|trial_flux_He|; columns \verb|F_H|,
   \verb|F_He|
 & the header pair by \path{element_flux_closure.py}, through
   \path{EXHALE_resolved.out}. The \verb|F_H|/\verb|F_He| \emph{columns} are
   written by the adapters and read by no Fortran routine: EXHALE measures
   its own elemental fluxes instead and reports them, which is the half of
   the loop it owns \\[2pt]
molecular/atomic H partition
 & columns \verb|q_H2|, \verb|q_H|
 & \verb|q_H2| $\to$ \verb|q_h2_base|, the base particle count.
   \verb|q_H| is required of the file and read by nothing \\[2pt]
$K_{zz}$ at the match
 & column \verb|Kzz|
 & not the match alone: the whole column, onto \verb|kzz_cell|
   (\verb|eddy_diffusion_on_grid|) \\[2pt]
selected molecular abundances the EXHALE network needs
 & \verb|q_H2|, plus the diagnostic set (\verb|q_H2O|, \verb|q_CO|,
   \verb|q_CH4|, \dots)
 & \emph{not met.} The EXHALE network resolves H$_2$, H$_2^+$, H$_3^+$ and
   HeH$^+$; the file carries none of the three ions and no molecular
   profile is seeded from it. The diagnostic columns are metadata
   (nothing consumes them), and the schema says so \\[2pt]
source fingerprint
 & header \verb|solution_id|, \verb|source_code|, \verb|source_version|,
   \verb|mechanism|, \verb|stellar_flux|, \verb|iteration|, \verb|notes|
 & the reader, and \verb|lap_report_provenance| into
   \path{EXHALE_resolved.out} \\
\bottomrule
\end{tabular}
\end{center}

The two further requirements the plan states are both met:

\begin{itemize}
\item \textbf{the same-solution consistency of the handoff pair is
  enforced}, not printed. With a profile in use a \path{base.inp} without a
  \verb|# solution_id| line, or with one that differs from the profile's,
  stops the run, and a scalar physics key beside the profile stops it
  earlier still (\path{input_read.f90:1344-1367,1371-1386}). Note the scope:
  this is enforcement of the \emph{pair}. A \path{base.inp} on its own is
  still checked only by the two printed messages of
  section~\ref{sec:scalars};
\item \textbf{the transit tool reads the resolved configuration.}
  \verb|read_input_parameters| prefers \path{EXHALE_resolved.out} next to
  \path{input.inp} and takes the planet radius, mass, equilibrium
  temperature, orbital distance, stellar mass and He/H from it, refusing to
  fall back silently if the file is present but unreadable
  (\path{exhale_transit_lib.py:336-364}). \path{EXHALE_resolved.out} is
  written after both handoff readers have run
  (\path{src/modules/files_IO/write_setup_report.f90:428-543}), so a handoff
  that moves the base radius now moves the transit geometry with it,
  Phase\,B of the plan.
\end{itemize}

\subsection{What is still open}

\begin{itemize}
\item The molecular ions of the network are not part of the handoff, as
  above. A photochemical H$_2$/H$_2^+$/H$_3^+$/HeH$^+$ solution at the base
  is still something EXHALE re-derives rather than inherits.
\item The overlap window the design defines the closure on was measured on
  LHS\,1140\,b and rejected there: it is $\sim\!0.003\,R_p$ thick and the
  flux across it spreads by a factor of 13--35, so the driver falls back on
  the steady-flux window $r \ge r_{\rm esc}$ and records that it did
  (\path{src/utils/element_flux_closure.py}, header). The fallback is
  reported rather than silent, but the window the design was written around
  is not the window the number comes from.
\item \verb|q_H2_base| and the profile's \verb|q_H2| are the composition of
  the inflowing gas in both routes, no longer EOS anchors alone (\S117 of
  \path{docs/Update_EXHALE_stage1.tex}, 2026-09-01). They set the base particle
  count \emph{and} are imposed on the H$_2$ partition of the lower ghost
  species, the particle count then coming from that same species state.
  Before that change the two disagreed by 6.4\% and the base ran 73~K above
  the $T_0$ the boundary condition claimed to impose.
\end{itemize}

\end{document}
